\subsection{PROPERTIES OF UNIFORMITIES DEFINED BY FAMILIES OF PSEUDOMETRICS}

Let \(\mathcal{U}\) be a uniformity defined on a set \(X\) by a family of finite pseudometrics \((f_{i})\) . If we endow \(X \times X\) with the uniformity which is the product of \(\mathcal{U}\) by itself, then each of the real-valued functions \(f_{i}\) is uniformly continuous on \(X \times X\) ; for by (1) we have

\[
\left| f _ {t} (x, y) - f _ {t} \left(x ^ {\prime}, y ^ {\prime}\right) \right| \leqslant f _ {t} (x, x ^ {\prime}) + f _ {t} (y, y ^ {\prime}),
\]

and therefore the relations \(f_{i}(x, x') \leqslant \varepsilon / 2\) , \(f_{i}(y, y') \leqslant \varepsilon / 2\) imply

\[
\left| f _ {t} (x, y) - f _ {t} \left(x ^ {\prime}, y ^ {\prime}\right) \right| \leqslant \varepsilon .
\]

For \(\mathcal{U}\) to be Hausdorff it is necessary and sufficient, from the definition of the entourages of \(\mathcal{U}\) , that for each pair of distinct points \(x, y\) of \(\mathbf{X}\) there is an index \(\iota\) such that \(f_{\iota}(x, y) \neq 0\) .

In particular, if \(\mathcal{U}\) is defined by a single pseudometric \(f\) , then \(\mathcal{U}\) is Hausdorff if and only if the relation \(f(x,y)=0\) implies \(x=y\) (cf.~§ 2). If \(\mathcal{U}\) is not Hausdorff, the intersection of all the entourages of \(\mathcal{U}\) is the subset of \(\mathbf{X} \times \mathbf{X}\) consisting of pairs \((x,y)\) such that \(f_{i}(x,y)=0\) for all \(\iota\) ; this subset is the graph of an equivalence relation \(\mathbf{R}\) on \(\mathbf{X}\) , and the Hausdorff uniformity associated with \(\mathcal{U}\) is defined on \(\mathbf{X}/\mathbf{R}\) (cf.~Chapter II, § 3, no. 8). It is then easily verified that the functions \(f_{i}\) are compatible (in \(x\) and in \(y\) ) with the relation \(\mathbf{R}\) (Set Theory, \(\mathbf{R}\) , § 5, no. 7) and that the functions \(\overline{f_{i}}\) , obtained from \(f_{i}\) by passing to the quotient (with respect to \(x\) and \(y\) ), are pseudometrics on \(\mathbf{X}/\mathbf{R}\) which define the Hausdorff uniformity associated with \(\mathcal{U}\) (cf.~§ 2, no. 1).

If A is a non-empty subset of X, the restriction to \(A \times A\) of a pseudometric on X is clearly a pseudometric on A. The uniformity induced by U on A is clearly that defined by the family of restrictions to \(A \times A\) of the pseudometrics \(f_{t}\) .

Let us now look at the completion of the uniform space \(\mathbf{X}\) when \(\mathcal{U}\) is Hausdorff.

PROPOSITION 1. Let X be a Hausdorff uniform space whose uniformity U is defined by a family of finite pseudometrics ( \(f_{t}\) ), and let \(\hat{X}\) be the completion of X. Then the functions \(f_{t}\) can be extended by continuity to \(\hat{\mathbf{X}} \times \hat{\mathbf{X}}\) ; the extended functions \(\overline{f_t}\) are finite pseudometrics on \(\hat{\mathbf{X}} \times \hat{\mathbf{X}}\) , and the family \((f_{t})\) defines the uniformity of \(\hat{\mathbf{X}}\) .

First, the \(f_{t}\) can be extended by continuity to \(\hat{\mathbf{X}}\times\hat{\mathbf{X}}\) , because they are uniformly continuous on \(\mathbf{X}\times\mathbf{X}\) ; and the extended functions \(\overline{f_{t}}\) are uniformly continuous on \(\hat{\mathbf{X}}\times\hat{\mathbf{X}}\) (Chapter II, § 3, no. 6, Theorem 2); moreover, they are pseudometrics on \(\hat{\mathbf{X}}\) by virtue of the principle of extension of inequalities (Chapter IV, § 5, no. 2, Theorem 1). Let \(\mathcal{U}_{1}\) denote the uniformity on \(\hat{\mathbf{X}}\) obtained by completion, and let \(\mathcal{U}_{2}\) denote the uniformity defined by the family of pseudometrics ( \(\overline{f_{t}}\) ). Then \(\mathcal{U}_{2}\) is coarser than \(\mathcal{U}_{1}\) ; for each \(\overline{f_{t}}\) is uniformly continuous on \(\hat{\mathbf{X}}\times\hat{\mathbf{X}}\) with respect to \(\mathcal{U}_{1}\) , and hence for each \(a>0\) there exists an entourage V of \(\mathcal{U}_{1}\) such that, whenever \((x,y)\in\mathrm{V}\) , we have \(|\overline{f_{t}}(x,y)-\overline{f_{t}}(x,x)|\leqslant a\) , that is {[}since \(\overline{f_{t}}(x,x)=0\) {]}, \(\mathrm{V}\subset\overline{f_{t}}([0,a])\) ; hence every entourage of \(\mathcal{U}_{2}\) is an entourage of \(\mathcal{U}_{1}\) . On the other hand, \(\mathcal{U}_{1}\) and \(\mathcal{U}_{2}\) induce the same uniformity \(\mathcal{U}\) on X. As \(\hat{\mathbf{X}}\) is complete with respect to \(\mathcal{U}_{1}\) , it follows that \(\mathcal{U}_{1}\) and \(\mathcal{U}_{2}\) are identical (Chapter II, § 3, no. 7, Proposition 14).

